% Sección sobre objetivos y control en seguridad radiológica
\section*{Objetivos}
Mantener controlada la dosis de personal. \\
Nivel de Investigación: \textbf{400 $\text{mRem}$ acumulados al mes} \\
Nivel de Intervención: \textbf{600 $\text{mRem}$ acumulados al mes}

Para rebasar algún límite se debe llenar la solicitud SR-03: Autorización de Dosis Adicional.

\section*{Control de fuentes radiactivas}
\subsection*{Solo Personal administrativo}

\section*{Almacenamiento de Fuentes Radiactivas}
El almacén debe cumplir con:
\begin{itemize}
    \item Tener acceso controlado con llave.
    \item Al resguardo de condiciones ambientales.
    \item Ser instalación fija, no portátil.
    \item Nivel de Radiación en el exterior no mayor a 5 mR/hr.
    \item Estar señalizado.
    \item No almacenar material explosivo o flamable junto.
\end{itemize}

\section*{Transporte de Fuentes Radiactivas}
El vehículo contará con:
\begin{itemize}
    \item Calcomanías de radiación.
    \item Extintor.
    \item Juego de herramientas.
    \item Caja portacontenedor con candado.
    \item Debe estar en buenas condiciones de circulación.
\end{itemize}
Las fuentes no deben transportarse en compartimentos ocupados por pasajeros (autobús, ferrocarril, avión, correo, etc.).

\section*{Control de material y equipo de Seguridad Radiológica}
El equipo se controla por residencias y con el uso de los formatos AL-01 "RESGUARDO DE MATERIAL Y EQUIPO" y AL-02 "TRANSFERENCIA DE MATERIAL Y EQUIPO".

\section*{Protección Radiológica en la toma de radiografías}
\begin{itemize}
    \item Verificar los dosímetros.
    \item Verificar la alarma (con el contenedor).
    \item Verificar Detector (Baterías, vigencia de calibración y respuesta)
    \item Verificar equipo de gammagrafía: No golpes, no estrangulamientos.
    \item Acordonar: 15 m si es gamma, 10 m si es rayos x.
    \item Señalizar: Distancias no mayores de 10 m entre signos.
    \item Mientras se esté exponiendo hay que vigilar el área.
\end{itemize}
